\section{Proof of the classification}\label{sec:classification}

\subsection{Necessity}
Suppose the circle set in Theorem~\ref{thm:main} is Ramsey.
Proposition~\ref{prop:necessity} gives a tensor certificate.
Transport and descent give a certificate for the finite rational-frame
normalization of Section~\ref{sec:second-order}. In its moduli field,
$\dim V=3$ and $V\cap W=0$. Proposition~\ref{prop:graph} provides
a three-dimensional space $U$ satisfying the product property,
with $V\subset U+W$. Proposition~\ref{prop:gluing} gives $U=EW$,
and Lemma~\ref{lem:sign-groups} partitions the nodes into two
groups of size at least three, on each of which every derivation
has quadratic parameter velocity.

For four distinct finite parameters in one group put
\[
 \rho=\frac{(t_a-t_c)(t_b-t_d)}{(t_a-t_d)(t_b-t_c)}.
\]
If $v_i=D(t_i)$, the quotient rule gives
\begin{equation}\label{eq:cross-derivative}
 \frac{D\rho}{\rho}=
 \frac{v_a-v_c}{t_a-t_c}+\frac{v_b-v_d}{t_b-t_d}
 -\frac{v_a-v_d}{t_a-t_d}-\frac{v_b-v_c}{t_b-t_c}.
\end{equation}
For quadratic velocities $v_i=a+bt_i+ct_i^2$, each difference
quotient is $b+c(t_i+t_j)$, so the four terms cancel. Thus every
$\Q$-derivation of $F$ kills $\rho$. The constants characterization
in Section~\ref{sec:second-order} makes $\rho$ algebraic over
$\Q$. All internal cross ratios in both groups are algebraic.
This proves the necessary implication using neither direction of
OpenAI's Theorem~1.1.

\subsection{The known sufficient mechanism}\label{sec:positive}
For completeness we record the circle certificate behind the
sufficient implication. This is a case of Strinz's two-group
theorem~\cite[Theorem C2]{strinz2026spherical}; the conclusion
that the certificate implies Ramsey is OpenAI's sufficient direction.

In the unit circle coordinate field $F$, choose
$A_I,A_J\in\mathrm{GL}_2(F)$ taking three points of each group to
$0,1,\infty$. The other normalized parameters are real algebraic
elements of $F$. Choose algebraic homogeneous representatives
$w_i$, so $A_Hz_i=\lambda_iw_i$ with $\lambda_i\in F^\times$
for $i\in H$, $H=I,J$.

There is an element $e\in\cB_F$ such that
\begin{equation}\label{eq:algebraic-projector}
 m(e)=1,\qquad e(L(\alpha)-R(\alpha))=0
\end{equation}
for every coordinate $\alpha$ of the finitely many $w_i$.
Here is an explicit construction. If $f_\alpha\in\Q[X]$ is the
minimal polynomial, put
\[
 g_\alpha(X)=\frac{f_\alpha(X)}{X-\alpha}\in F[X],\qquad
 e_\alpha=L(f_\alpha'(\alpha)^{-1})
                       g_\alpha^L(R(\alpha)),\qquad
 e=\prod_\alpha e_\alpha.
\]
The superscript $L$ embeds the coefficients in the left factor.
The identities
$(R(\alpha)-L(\alpha))g_\alpha^L(R(\alpha))=0$
and $g_\alpha(\alpha)=f_\alpha'(\alpha)\ne0$ prove
\eqref{eq:algebraic-projector}.

For $z=(s,t)^T$, $w=(u,v)^T$, define the bilinear graph forms
\[
 K_H(z,w)=\det\bigl(L(A_H)z,R(A_H)w\bigr)\qquad(H=I,J).
\]
At a node of group $H$, \eqref{eq:algebraic-projector} gives
$eK_H(L(z_i),R(z_i))=0$. Multiplication of coefficients gives
$m(K_H)=\det(A_H)(sv-tu)$. Hence
\begin{equation}\label{eq:two-group-certificate}
 G(z,w)=-2L((\det A_I\det A_J)^{-1})eK_I(z,w)K_J(z,w)
\end{equation}
vanishes at every tensor node, and $m(G)=-2(sv-tu)^2$.
The basis identity \eqref{eq:circle-kernel} yields a matrix with
$m(P)=\HH$. Dividing the node vectors by $s_i^2+t_i^2>0$
gives \eqref{eq:tensor-criterion}. All denominators are nonzero
elements of $F$ in pure tensor factors; no domain assumption is used.
At least three distinct circle points span the plane, so OpenAI's
Theorem~1.1, sufficient direction, applies. This proves
Theorem~\ref{thm:main} and hence Corollary~\ref{cor:seven}.

In particular, the construction applies to a group of four points
with algebraic cross ratio and any remaining triple, without a
dimension restriction. More generally it includes the algebraic-plus-
three conclusion of Strinz. Both are positive Ramsey consequences
of the sufficient direction of OpenAI's Theorem~1.1. The scripts
\ancfile{structure.py} and \ancfile{eight_factory.py} check the
diagonal and graph identities used by these constructions.
